% 第五章：Linux内核实现

\newpage
\section{Linux内核实现} \label{sec:kernel}

本章详细介绍Yat-CASched算法在Linux内核中的具体实现过程。基于第四章提出的简化缓存感知策略，我们开发了一个完整的内核调度类，将理论算法转化为可在生产环境中运行的系统组件。

\subsection{实现概述}

基于简化的缓存感知策略，我们开发了一个名为\texttt{Yat\_Casched}的Linux内核调度类。该实现将简单缓存感知算法直接集成到操作系统内核中，专注于保持任务在其上次运行的CPU上执行。

\begin{tcolorbox}[
    colback=blue!5!white,
    colframe=blue!50!black,
    title=\textbf{内核实现核心目标},
    fonttitle=\bfseries,
    arc=3pt
]
\textbf{实现目标}：将理论算法转化为可在生产环境中运行的内核模块\\
\textbf{设计原则}：轻量级、高效率、兼容性好\\
\textbf{技术路径}：基于Linux CFS调度框架，最小化修改，最大化效果
\end{tcolorbox}

\subsection{调度策略集成}

我们在Linux内核中新增了调度策略\texttt{SCHED\_YAT\_CASCHED = 8}，作为非实时调度器的一种扩展：

\begin{tcolorbox}[
    colback=lightgreen!10!white,
    colframe=green!50!black,
    title=\textbf{调度策略参数配置},
    fonttitle=\bfseries,
    arc=3pt
]
\begin{itemize}
    \item \textbf{调度策略ID}：\texttt{SCHED\_YAT\_CASCHED = 8}
    \item \textbf{优先级设置}：非实时调度器，使用优先级0
    \item \textbf{CPU亲和性}：倾向于保持任务在上次运行的CPU上
    \item \textbf{设计理念}：轻量级实现，最小化调度开销
\end{itemize}
\end{tcolorbox}

\subsection{内核数据结构设计}

为支持Yat-CASched调度算法，我们在内核中定义了专门的数据结构：

\begin{tcolorbox}[
    colback=yellow!5!white,
    colframe=orange!50!black,
    title=\textbf{核心数据结构实现},
    fonttitle=\bfseries,
    arc=3pt
]
\begin{code}
// 每个任务的调度实体结构
struct sched_yat_casched_entity {
    u64 slice;           // 时间片长度
    int last_cpu;        // 上次运行的CPU编号
};

// 每个CPU的运行队列结构  
struct yat_casched_rq {
    struct task_struct *agent;  // 当前代理任务指针
};
\end{code}
\end{tcolorbox}

这些结构体被集成到Linux内核的任务调度框架中，通过条件编译\texttt{CONFIG\_SCHED\_CLASS\_YAT\_CASCHED}进行控制。

% 图片预留位置：内核数据结构关系图
\begin{figure}[htbp]
\centering
% \includegraphics[width=0.85\textwidth]{img/kernel_data_structure.png}
\begin{tcolorbox}[
    colback=gray!10,
    colframe=gray!50,
    width=0.85\textwidth,
    center,
    arc=3pt
]
\centering
\textbf{[图片预留位置]}\\
\textbf{内核数据结构关系图}\\
\footnotesize{展示调度实体、运行队列与Linux调度框架的集成关系}
\end{tcolorbox}
\caption{内核数据结构关系}
\label{fig:kernel_structure}
\end{figure}

\subsection{调度类接口实现}

Yat-CASched调度类实现了Linux内核调度框架要求的所有接口函数：

\begin{tcolorbox}[
    colback=lightblue!5!white,
    colframe=blue!50!black,
    title=\textbf{调度类完整接口定义},
    fonttitle=\bfseries,
    arc=3pt
]
\begin{code}
// 调度类定义 - Linux内核调度接口
DEFINE_SCHED_CLASS(yat_casched) = {
    .enqueue_task       = enqueue_task_yat_casched,    // 任务入队
    .dequeue_task       = dequeue_task_yat_casched,    // 任务出队
    .yield_task         = yield_task_yat_casched,      // 任务让出
    .yield_to_task      = yield_to_task_yat_casched,   // 让出给指定任务
    .wakeup_preempt     = wakeup_preempt_yat_casched,  // 唤醒抢占
    .pick_next_task     = pick_next_task_yat_casched,  // 选择下一个任务
    .put_prev_task      = put_prev_task_yat_casched,   // 放置前一个任务
    .set_next_task      = set_next_task_yat_casched,   // 设置下一个任务
#ifdef CONFIG_SMP
    .balance            = balance_yat_casched,         // SMP负载均衡
    .select_task_rq     = select_task_rq_yat_casched,  // 选择运行队列
    .set_cpus_allowed   = set_cpus_allowed_yat_casched,// 设置CPU亲和性
    .rq_online          = rq_online_yat_casched,       // 队列上线
    .rq_offline         = rq_offline_yat_casched,      // 队列下线
#endif
    .pick_task          = pick_task_yat_casched,       // 选择任务
    .task_tick          = task_tick_yat_casched,       // 任务时钟滴答
    .switched_to        = switched_to_yat_casched,     // 切换到调度器
    .prio_changed       = prio_changed_yat_casched,    // 优先级变更
    .update_curr        = update_curr_yat_casched,     // 更新当前任务
};
\end{code}
\end{tcolorbox}

\subsection{核心分配策略实现}

实现的核心分配机制包括：

\begin{tcolorbox}[
    colback=green!5!white,
    colframe=green!50!black,
    title=\textbf{CPU亲和性优化实现},
    fonttitle=\bfseries,
    arc=3pt
]
\textbf{1. CPU亲和性优化}
\begin{code}
int select_task_rq_yat_casched(struct task_struct *p,
                               int task_cpu, int flags) {
    // 核心策略：优先选择任务上次运行的CPU
    if (p->yat_casched.last_cpu >= 0 && 
        cpu_online(p->yat_casched.last_cpu)) {
        return p->yat_casched.last_cpu;  // 优先选择上次CPU
    }
    return task_cpu;  // 回退到默认CPU分配
}
\end{code}
\end{tcolorbox}

\begin{tcolorbox}[
    colback=orange!5!white,
    colframe=orange!50!black,
    title=\textbf{任务切换处理实现},
    fonttitle=\bfseries,
    arc=3pt
]
\textbf{2. 任务切换处理}
\begin{code}
void put_prev_task_yat_casched(struct rq *rq,
                               struct task_struct *p) {
    // 关键操作：记录任务上次运行的CPU编号
    // 为下次调度提供缓存感知依据
    p->yat_casched.last_cpu = rq->cpu;
}
\end{code}
\end{tcolorbox}

% 图片预留位置：CPU亲和性调度过程图
\begin{figure}[htbp]
\centering
% \includegraphics[width=0.9\textwidth]{img/cpu_affinity_process.png}
\begin{tcolorbox}[
    colback=gray!10,
    colframe=gray!50,
    width=0.9\textwidth,
    center,
    arc=3pt
]
\centering
\textbf{[图片预留位置]}\\
\textbf{CPU亲和性调度过程示意图}\\
\footnotesize{展示任务在多核间的调度决策过程和缓存感知优化}
\end{tcolorbox}
\caption{CPU亲和性调度过程}
\label{fig:cpu_affinity}
\end{figure}

\subsection{系统集成与配置}

\subsubsection{内核配置选项}

在内核配置系统中添加了专门的配置选项：

\begin{code}
# kernel/sched/Kconfig
config SCHED_CLASS_YAT_CASCHED
    bool "Yat_Casched scheduling class"
    default y
    help
      启用Yat_Casched调度类。该调度器通过优先保持任务在
      上次使用的CPU上运行来提供CPU亲和性优化。
\end{code}

\subsubsection{编译系统集成}

通过修改内核Makefile实现模块化编译：

\begin{code}
# kernel/sched/Makefile
obj-$(CONFIG_SCHED_CLASS_YAT_CASCHED) += yat_casched.o
\end{code}

\subsubsection{用户态接口}

为用户程序提供了标准的调度策略设置接口：

\begin{code}
// 设置进程使用Yat-CASched调度策略
struct sched_param param;
param.sched_priority = 0;  // 非实时调度器，优先级为0

int ret = sched_setscheduler(pid, SCHED_YAT_CASCHED, &param);
if (ret == 0) {
    printf("Successfully set Yat_Casched policy\n");
}
\end{code}

\subsection{功能验证与基础测试}

\subsubsection{测试程序设计}

我们开发了综合测试程序\texttt{test\_yat\_casched.c}来验证调度器的功能：

\begin{itemize}
    \item \textbf{调度策略切换测试}：验证进程/线程能否成功切换到Yat-CASched策略
    \item \textbf{多线程CPU亲和性测试}：监控线程在不同CPU之间的切换情况
    \item \textbf{调度行为验证}：记录CPU切换次数，验证缓存局部性优化效果
\end{itemize}

\subsubsection{内核编译与集成}

实现了完整的内核集成流程：

\textbf{1. 内核配置}
\begin{code}
# 启用Yat-CASched调度器
make menuconfig
# 选择：Processor type and features -> 
#       Yat_Casched scheduling class
\end{code}

\textbf{2. 编译验证}
\begin{code}
# 编译内核
make -j$(nproc)

# 验证调度器集成
nm vmlinux | grep yat_casched
\end{code}

\textbf{3. 测试环境部署}
\begin{code}
# 创建测试initramfs
./create_simple_initramfs.sh

# 启动QEMU测试环境
qemu-system-x86_64 -kernel arch/x86/boot/bzImage \
    -initrd initramfs.cpio.gz -m 2048 -smp 4 \
    -append "console=ttyS0 init=/init"
\end{code}

\subsubsection{基础功能验证结果}

通过在QEMU虚拟环境中的测试，验证了以下关键功能：

\textbf{1. 调度策略注册成功}
\begin{itemize}
    \item 系统成功识别\texttt{SCHED\_YAT\_CASCHED = 8}调度策略
    \item 用户程序能够通过\texttt{sched\_setscheduler()}切换到新调度策略
    \item 调度器在\texttt{/proc/sched\_debug}中正确显示
\end{itemize}

\textbf{2. CPU亲和性优化效果}
\begin{itemize}
    \item 任务倾向于在同一CPU上连续执行
    \item 相比\texttt{SCHED\_OTHER}，CPU切换频率显著降低
    \item 多线程环境下保持良好的负载均衡
\end{itemize}

\textbf{3. 系统稳定性验证}
\begin{itemize}
    \item 调度器在多核环境下稳定运行
    \item 与其他调度策略兼容，无冲突
    \item 内核模块可正常加载和卸载
\end{itemize}

\begin{tcolorbox}[
    colback=blue!10!white,
    colframe=blue!75!black,
    title=\textbf{本章技术贡献总结},
    fonttitle=\bfseries,
    arc=3pt
]
\begin{itemize}
    \item \textbf{系统集成}：成功将算法集成到Linux内核，实现了\texttt{SCHED\_YAT\_CASCHED}调度类
    \item \textbf{数据结构设计}：定义了简洁高效的调度实体和运行队列结构
    \item \textbf{接口实现}：完成了Linux调度框架要求的所有核心接口函数
    \item \textbf{配置系统}：建立了完整的内核配置、编译和部署流程
    \item \textbf{基础验证}：验证了调度器的基本功能和系统稳定性
\end{itemize}
\end{tcolorbox}

\begin{tcolorbox}[
    colback=green!5!white,
    colframe=green!50!black,
    title=\textbf{面向第六章的测试准备},
    fonttitle=\bfseries,
    arc=3pt
]
本章的内核实现为第六章的详细性能测试奠定了坚实基础：
\begin{itemize}
    \item \textbf{实验环境}：QEMU/KVM虚拟化平台已就绪
    \item \textbf{测试工具}：内核模块和用户态测试程序开发完成
    \item \textbf{监控机制}：调度器状态监控和数据收集框架已建立
    \item \textbf{对比基准}：与传统CFS调度器的对比测试环境已配置
\end{itemize}
第六章将展示详细的性能测试结果和可视化分析，量化验证缓存感知调度的实际效果。
\end{tcolorbox}
